\chapter{Benchmarking Cryptographic Primitives}

\section{Operation-Level Benchmarks}

Operation-level benchmarks isolate primitive kernels inside the zkVM and measure
their contribution to proving cost. The workload set focuses on operations that
are expected to dominate cryptographic traces, including substitution layers,
matrix-heavy transforms, and bitwise round-function kernels.

This focus is motivated by proof-system cost structure. While RISC Zero is a
VM-centric proving stack rather than a hand-authored arithmetic circuit, prover
work still depends on enforcing constraints over the execution trace. In
practice, operation shape and data movement patterns can increase trace pressure
and therefore proving time \cite{bensasson2018stark,risczeroTraceDocs}.

These measurements are used to rank optimization direction, then cross-checked
against full cipher runs. They are evidence for engineering priority, not
standalone performance claims about complete cryptographic pipelines.

\section{Cipher-Level Benchmarks}

Cipher-level benchmarks execute complete implementations under the same RISC Zero
proving flow. The core comparison is AES, Salsa20, and LowMC baselines, with
optimized variants analyzed separately.

Interpretation rule: zkVM proving cost does not necessarily mirror native CPU
performance intuition, so conclusions are based on measured prover behavior.
Where operation-level and cipher-level conclusions disagree, the cipher-level
result is treated as authoritative for system-design decisions.

\section{Cross-Algorithm Results}

Cross-algorithm comparison values are taken from the baseline comparison
snapshot used in this report.

This run provides a like-for-like single-trial snapshot containing
\texttt{aes-r0}, \texttt{salsa-r0}, \texttt{lowmc-r0}, and
\texttt{lowmc-r0-optimised}.

\begin{table}[t]
\centering
\small
\begin{tabular}{lrrrr}
\hline
Benchmark ID & Prove s (median) & Verify s (median) & User cycles (median) & Success \\
\hline
\texttt{aes-r0} & 14.843 & 0.0111 & 181710 & 1.0 \\
\texttt{salsa-r0} & 29.789 & 0.0117 & 467020 & 1.0 \\
\texttt{lowmc-r0} & 695.050 & 0.1330 & 10388894 & 1.0 \\
\texttt{lowmc-r0-optimised} & 517.209 & 0.1065 & 7919919 & 1.0 \\
\hline
\end{tabular}
\caption{Cross-algorithm comparison from the baseline run.}
\label{tab:cross-algo-comparison}
\end{table}

Even in this single-trial snapshot, the effect-size gap is large: LowMC remains
substantially slower than AES, while Salsa20 sits between the two.

The important interpretation is not merely that AES is faster in Table
\ref{tab:cross-algo-comparison}. The result shows that the naive ordering
suggested by circuit-friendliness is not preserved by this VM implementation:
LowMC has the strongest ZK-oriented design motivation, but the worst observed
proving time among the tested targets in this snapshot. Salsa20 behaves as a
useful middle point, suggesting that software regularity helps but is not enough
to beat the compact AES baseline in this specific stack.

This also explains how the result is used later. The table is not a universal
claim about primitive quality; it is a decision input for the implemented RISC
Zero pipeline. It justifies pivoting the final authenticated construction toward
AES while keeping LowMC analysis as evidence for why VM-level measurements are
needed.

\section{Operations Directional Ranking}

This directional ranking uses benchmark IDs
\texttt{operation-bnchmrk-r0:*} from the operations campaign. Trial
policy is one trial per operation row, and all listed rows succeeded
(\texttt{success\_rate = 1.0}).

The operation rows separate into coarse trace buckets. Addition and subtraction
fall into the 32768-cycle bucket and prove in about 1.8--1.9 seconds, while the
boolean, rotation, multiplication, and Toffoli-style rows fall into the
65536-cycle bucket and cluster around 3.6--3.7 seconds. Within that larger bucket,
rotations and multiplication have the highest user-cycle counts, so they are
useful hotspot indicators even though the rounded trace size hides some fine
ordering in wall-clock time.

\begin{figure}[t]
\centering
\includegraphics[width=0.95\linewidth]{Figures/report-benchmarks/operations_directional_prove_time.png}
\caption{Operation-level directional ranking by median proving time from the report operations campaign.}
\label{fig:operations-directional-prove}
\end{figure}

\section{Operation-to-Cipher Link Note}

The operation campaign and cipher campaign align on one key expectation:
expensive bitwise and nonlinear kernels are not enough to make a primitive
zkVM-efficient by themselves. Full-cipher composition overhead dominates.

\textbf{Matches expectation:} LowMC's matrix/bit-transformation-heavy structure
corresponds to very high full-cipher proving cost relative to AES.

\textbf{Does not match simple expectation:} no single operation row in
Figure~\ref{fig:operations-directional-prove} explains the full \(\sim 35\times\)
cipher-level gap between LowMC and AES seen later. The gap emerges from repeated
round composition, memory movement, and trace-length accumulation, not one
isolated kernel.

The key insight is therefore compositional: operation rows can identify plausible
pressure points, but full-cipher proving cost is dominated by how often those
operations are repeated and how much state they move through the guest program.
For that reason, operation-level evidence is useful for optimization direction
but insufficient for final primitive selection.
